/*
 * Automacao residencial - PIC18F4550 no kit XM118
 *
 * Ligacoes / dip switches:
 *   Botoes   : RB0, RB1, RB2 (SW4, SW5, SW6) - dips do PORTB (CH6) em OFF
 *   Luzes    : RD0, RD1, RD2 (LEDs do PORTD, ativos em nivel baixo no kit)
 *   LM35     : RA0/AN0   -> CH1 chave 7 (TEMP) ON
 *   Ventoinha: RC2/CCP1  -> CH3 chave 5 (VENTILADOR) ON (PWM)
 *   Tacografo: RC0/T1CKI -> CH3 chave 2 (TACOGRAFO) ON
 *   LCD      : conforme nxlcd.h -> CH2 chave 1 ON, bicolores (CH2 5..8) OFF
 */

#include <xc.h>
#include <stdint.h>
#include <stdio.h>
#include "nxlcd.h"

#pragma config CPUDIV = OSC1_PLL2
#pragma config FOSC = HS
#pragma config WDT = OFF
#pragma config PBADEN = OFF
#pragma config LVP = OFF
#pragma config DEBUG = OFF
#pragma config MCLRE = ON

#define _XTAL_FREQ 20000000

// ---------- Configuracoes ----------
#define BTN_RB0  0x01
#define BTN_RB1  0x02
#define BTN_RB2  0x04

#define NUM_AMB          3
#define NIVEL_ON         0      // kit XM118: LEDs ativos em nivel BAIXO (use 1 no Proteus se for ativo alto)

// Segundos de cada unidade configurada nos timers:
// 60 = minutos | 3600 = horas | 1 = segundos (bom para testar)
#define SEG_POR_UNIDADE  60UL
#define UNID_TXT         "min"
#define CONT_MAX         99

#define PULSOS_POR_VOLTA 7      // numero de pas da ventoinha (ajuste para calibrar o rpm)

#define MODO_AUTO        0
#define MODO_LIGADA      1
#define MODO_DESLIGADA   2

// Enderecos da EEPROM interna
#define EE_MAGIC   0
#define EE_LIMITE  1
#define EE_MODO    2
#define EE_VEL     3

// ---------- Variaveis globais ----------
volatile uint32_t t_on[NUM_AMB];     // segundos ate ligar  (0 = sem timer)
volatile uint32_t t_off[NUM_AMB];    // segundos ate desligar (0 = sem timer)
volatile uint8_t  luz[NUM_AMB];      // estado desejado de cada luz
volatile uint8_t  tick = 0;          // ISR avisa: passou 1 segundo
volatile uint16_t pulsos = 0;        // pulsos do tacografo no ultimo segundo
volatile uint8_t  evento_amb = 0;    // 0 = nada | n = ambiente (n-1) mudou pelo timer
volatile uint8_t  evento_est = 0;    // novo estado da luz do evento

uint16_t temp_dec = 0;               // temperatura em decimos de grau
uint16_t rpm_atual = 0;
uint8_t  fan_on = 0;

uint8_t  limite = 30;                // limite de temperatura (C)
uint8_t  modo = MODO_AUTO;
uint8_t  vel = 100;                  // velocidade da ventoinha (%)

char lcd_buf[17];
char num_buf[24];

const char *nome_amb[NUM_AMB] = {"Sala", "Quarto", "Cozinha"};
const char *nome_modo[3]      = {"AUTO", "LIGADA", "DESLIG."};
const char *itens_menu[5]     = {"Iluminacao", "Temperatura", "Ventilacao", "Timers", "Limite temp."};
#define NUM_ITENS 5

// ---------- Interrupcao: 1 segundo ----------
void __interrupt() ISR(void)
{
    if(INTCONbits.TMR0IF == 1)
    {
        uint8_t i;
        uint16_t p;

        TMR0H = 0x67;
        TMR0L = 0x69;
        INTCONbits.TMR0IF = 0;

        for(i = 0; i < NUM_AMB; i++)
        {
            if(t_on[i] > 0)
            {
                t_on[i]--;
                if(t_on[i] == 0)
                {
                    luz[i] = 1;
                    evento_amb = i + 1;
                    evento_est = 1;
                }
            }
            if(t_off[i] > 0)
            {
                t_off[i]--;
                if(t_off[i] == 0)
                {
                    luz[i] = 0;
                    evento_amb = i + 1;
                    evento_est = 0;
                }
            }
        }

        // Tacografo: le e zera o Timer1 (RD16 = 1: ler L primeiro, escrever H primeiro)
        p = TMR1L;
        p |= ((uint16_t)TMR1H << 8);
        TMR1H = 0;
        TMR1L = 0;
        pulsos = p;

        tick = 1;
    }
}

// ---------- EEPROM interna ----------
uint8_t eeprom_le(uint8_t end)
{
    EEADR = end;
    EECON1bits.EEPGD = 0;
    EECON1bits.CFGS = 0;
    EECON1bits.RD = 1;
    return EEDATA;
}

void eeprom_grava(uint8_t end, uint8_t dado)
{
    uint8_t gie = INTCONbits.GIE;

    EEADR = end;
    EEDATA = dado;
    EECON1bits.EEPGD = 0;
    EECON1bits.CFGS = 0;
    EECON1bits.WREN = 1;

    INTCONbits.GIE = 0;
    EECON2 = 0x55;
    EECON2 = 0xAA;
    EECON1bits.WR = 1;
    INTCONbits.GIE = gie;

    while(EECON1bits.WR);
    PIR2bits.EEIF = 0;
    EECON1bits.WREN = 0;
}

void salva_config(void)
{
    eeprom_grava(EE_LIMITE, limite);
    eeprom_grava(EE_MODO, modo);
    eeprom_grava(EE_VEL, vel);
    eeprom_grava(EE_MAGIC, 0xA5);
}

void carrega_config(void)
{
    if(eeprom_le(EE_MAGIC) == 0xA5)
    {
        limite = eeprom_le(EE_LIMITE);
        modo   = eeprom_le(EE_MODO);
        vel    = eeprom_le(EE_VEL);

        if(limite < 15 || limite > 60) limite = 30;
        if(modo > 2) modo = MODO_AUTO;
        if(vel < 25 || vel > 100) vel = 100;
    }
}

// ---------- LCD ----------
void lcd_linha(uint8_t linha, const char *txt)
{
    uint8_t i = 0;

    while(i < 16 && txt[i] != '\0')
    {
        lcd_buf[i] = txt[i];
        i++;
    }
    while(i < 16)
    {
        lcd_buf[i] = ' ';
        i++;
    }
    lcd_buf[16] = '\0';

    WriteCmdXLCD(linha == 0 ? 0x80 : 0xC0);
    putsXLCD(lcd_buf);
}

// Retorna 1 uma unica vez por clique (debounce + espera soltar)
uint8_t pressionado(uint8_t mascara)
{
    if((PORTB & mascara) == 0)
    {
        __delay_ms(30);
        if((PORTB & mascara) == 0)
        {
            while((PORTB & mascara) == 0);
            __delay_ms(30);
            return 1;
        }
    }
    return 0;
}

// ---------- Luzes (somente liga/desliga, sem PWM) ----------
void escreve_led(uint8_t i, uint8_t ligado)
{
    uint8_t nivel = ligado ? NIVEL_ON : !NIVEL_ON;

    switch(i)
    {
        case 0: LATDbits.LATD0 = nivel; break;
        case 1: LATDbits.LATD1 = nivel; break;
        case 2: LATDbits.LATD2 = nivel; break;
    }
}

void aplica_luzes(void)
{
    uint8_t i;
    for(i = 0; i < NUM_AMB; i++)
        escreve_led(i, luz[i]);
}

// ---------- Temperatura (LM35 em AN0) ----------
// Retorna a temperatura em decimos de grau Celsius (ex.: 275 = 27,5 C)
uint16_t le_temperatura(void)
{
    uint32_t soma = 0;
    uint8_t i;

    for(i = 0; i < 8; i++)
    {
        ADCON0bits.GO = 1;
        while(ADCON0bits.GO);
        soma += ADRES;
    }
    soma >>= 3;

    // mV = adc * 5000 / 1024 ; LM35 = 10 mV/C -> decimos de grau = mV
    return (uint16_t)((soma * 5000UL) / 1024UL);
}

// ---------- Ventoinha (PWM no CCP1 / RC2) ----------
void pwm_ventoinha(uint8_t pct)
{
    uint16_t d = ((uint16_t)pct * 5) / 2;   // 0..250 (PR2 = 249)
    CCP1CONbits.DC1B = 0;
    CCPR1L = (uint8_t)d;
}

void controle_ventoinha(void)
{
    uint16_t lim = (uint16_t)limite * 10;

    if(modo == MODO_AUTO)
    {
        if(temp_dec >= lim)
            fan_on = 1;
        else if(temp_dec + 10 <= lim)       // histerese de 1 C
            fan_on = 0;
    }
    else
    {
        fan_on = (modo == MODO_LIGADA);
    }

    pwm_ventoinha(fan_on ? vel : 0);
}

// ---------- Tarefas de fundo (chamadas em todos os loops) ----------
void tarefas(void)
{
    uint16_t p;

    aplica_luzes();

    if(tick)
    {
        tick = 0;

        di();
        p = pulsos;
        ei();
        rpm_atual = (uint16_t)(((uint32_t)p * 60UL) / PULSOS_POR_VOLTA);

        temp_dec = le_temperatura();
        controle_ventoinha();
    }
}

// ---------- Menu ----------
void desenha_menu(int menu)
{
    char tmp[20];
    uint8_t primeiro = (uint8_t)menu & 0xFE;

    sprintf(tmp, "%c %s", (menu == primeiro) ? '>' : ' ', itens_menu[primeiro]);
    lcd_linha(0, tmp);

    if(primeiro + 1 < NUM_ITENS)
    {
        sprintf(tmp, "%c %s", (menu == primeiro + 1) ? '>' : ' ', itens_menu[primeiro + 1]);
        lcd_linha(1, tmp);
    }
    else
    {
        lcd_linha(1, "");
    }
}

// Ajuste generico de um valor: RB0 diminui, RB1 aumenta, RB2 confirma
void ajusta_valor(const char *titulo, uint8_t *v, uint8_t vmin, uint8_t vmax, const char *unid)
{
    lcd_linha(0, titulo);

    while(1)
    {
        tarefas();

        sprintf(num_buf, "%u %s", *v, unid);
        lcd_linha(1, num_buf);

        if(pressionado(BTN_RB0))
        {
            if(*v > vmin) (*v)--;
        }
        if(pressionado(BTN_RB1))
        {
            if(*v < vmax) (*v)++;
        }
        if(pressionado(BTN_RB2))
            return;
    }
}

// ---------- Telas ----------

void tela_luz(void)
{
    uint8_t a = 0;

    lcd_linha(1, "Prox   Alt   Sai");

    while(1)
    {
        tarefas();

        sprintf(num_buf, "%s: %s", nome_amb[a], luz[a] ? "lig." : "desl.");
        lcd_linha(0, num_buf);

        if(pressionado(BTN_RB0))
        {
            a++;
            if(a >= NUM_AMB) a = 0;
        }

        if(pressionado(BTN_RB1))
            luz[a] = !luz[a];

        if(pressionado(BTN_RB2))
            return;
    }
}

void tela_temperatura(void)
{
    uint16_t ultimo = 0xFFFF;

    sprintf(num_buf, "Limite: %u C", limite);
    lcd_linha(1, num_buf);

    while(!pressionado(BTN_RB2))
    {
        tarefas();

        if(temp_dec != ultimo)
        {
            ultimo = temp_dec;
            sprintf(num_buf, "Temp: %u.%u C", temp_dec / 10, temp_dec % 10);
            lcd_linha(0, num_buf);
        }
    }
}

void tela_ventilacao(void)
{
    while(1)
    {
        tarefas();

        sprintf(num_buf, "%s %u%%", nome_modo[modo], vel);
        lcd_linha(0, num_buf);

        sprintf(num_buf, "%u rpm %s", rpm_atual, fan_on ? "ON" : "OFF");
        lcd_linha(1, num_buf);

        if(pressionado(BTN_RB0))        // muda o modo
        {
            modo++;
            if(modo > 2) modo = 0;
            controle_ventoinha();
            salva_config();
        }

        if(pressionado(BTN_RB1))        // muda a velocidade
        {
            if(vel >= 100) vel = 25;
            else vel += 25;
            controle_ventoinha();
            salva_config();
        }

        if(pressionado(BTN_RB2))
            return;
    }
}

void tela_timer(void)
{
    uint8_t a = 0;
    uint8_t on_v = 0, off_v = 0;

    // 1) escolhe o ambiente
    lcd_linha(0, "Prox    OK   Sai");
    while(1)
    {
        tarefas();
        lcd_linha(1, nome_amb[a]);

        if(pressionado(BTN_RB0))
        {
            a++;
            if(a >= NUM_AMB) a = 0;
        }
        if(pressionado(BTN_RB1))
            break;
        if(pressionado(BTN_RB2))
            return;
    }

    // 2) tempos (0 = sem timer)
    ajusta_valor("Ligar apos:", &on_v, 0, CONT_MAX, UNID_TXT);
    ajusta_valor("Desligar apos:", &off_v, 0, CONT_MAX, UNID_TXT);

    // Desligar precisa vir depois de ligar
    if(on_v > 0 && off_v > 0 && off_v <= on_v)
    {
        lcd_linha(0, "Erro: desligar");
        lcd_linha(1, "deve ser > ligar");
        __delay_ms(1500);
        return;
    }

    di();
    t_on[a]  = (uint32_t)on_v  * SEG_POR_UNIDADE;
    t_off[a] = (uint32_t)off_v * SEG_POR_UNIDADE;
    ei();

    lcd_linha(0, nome_amb[a]);
    lcd_linha(1, (on_v > 0 || off_v > 0) ? "Timer salvo" : "Timer cancelado");
    __delay_ms(1000);
}

void tela_limite(void)
{
    ajusta_valor("Limite temp:", &limite, 15, 60, "C");
    salva_config();
    controle_ventoinha();

    lcd_linha(0, "Limite salvo");
    lcd_linha(1, "");
    __delay_ms(800);
}

// ---------- Programa principal ----------
void main(void)
{
    int menu = 0;
    uint8_t i;

    // Portas
    TRISA = 0xFF;               // RA0 = LM35
    TRISB = 0xFF;               // RB0..RB2 = botoes
    TRISC = 0b00000001;         // RC0 = tacografo (entrada), RC2 = PWM (saida)
    TRISD = 0x00;
    TRISE = 0x00;

    INTCON2bits.RBPU = 0;       // pull-ups internos do PORTB

    // ADC: somente AN0 analogico, Vref = VDD/VSS, resultado justificado a direita
    ADCON1 = 0b00001110;
    ADCON2 = 0b10101110;        // ADFM=1, Tacq = 12 TAD, Fosc/64
    ADCON0 = 0b00000001;        // canal AN0, ADC ligado

    // PWM da ventoinha (CCP1 em RC2): ~1,25 kHz
    PR2 = 249;
    CCPR1L = 0;
    CCP1CON = 0b00001100;
    T2CON = 0b00000111;         // Timer2 ligado, prescaler 1:16

    // Timer1 como contador de pulsos do tacografo (RC0/T1CKI)
    T1CON = 0b10000111;         // RD16=1, prescaler 1:1, clock externo, sem sincronismo, ligado
    TMR1H = 0;
    TMR1L = 0;

    // Estado inicial
    for(i = 0; i < NUM_AMB; i++)
    {
        luz[i] = 0;
        t_on[i] = 0;
        t_off[i] = 0;
    }
    aplica_luzes();
    carrega_config();

    // Timer0: 16 bits, Fosc/4, prescaler 1:128 -> estouro a cada ~1 s
    T0CONbits.T08BIT = 0;
    T0CONbits.T0CS = 0;
    T0CONbits.PSA = 0;
    T0CONbits.T0PS = 0b110;
    TMR0H = 0x67;
    TMR0L = 0x69;

    INTCONbits.TMR0IF = 0;
    INTCONbits.TMR0IE = 1;
    INTCONbits.PEIE = 1;
    INTCONbits.GIE = 1;

    T0CONbits.TMR0ON = 1;

    // Primeira leitura de temperatura
    temp_dec = le_temperatura();
    controle_ventoinha();

    // LCD
    OpenXLCD(FOUR_BIT & LINES_5X7);
    WriteCmdXLCD(0x01);
    __delay_ms(2);
    WriteCmdXLCD(0x0C);

    lcd_linha(0, " Menu da casa");
    lcd_linha(1, "   Bem vindo");
    __delay_ms(1000);

    desenha_menu(menu);

    while(1)
    {
        tarefas();

        if(evento_amb)
        {
            uint8_t a = evento_amb - 1;
            uint8_t e = evento_est;
            evento_amb = 0;

            lcd_linha(0, nome_amb[a]);
            lcd_linha(1, e ? "ligada p/ timer" : "desl. p/ timer");
            __delay_ms(1500);
            desenha_menu(menu);
        }

        if(pressionado(BTN_RB0))
        {
            menu++;
            if(menu >= NUM_ITENS)
                menu = 0;
            desenha_menu(menu);
        }

        if(pressionado(BTN_RB1))
        {
            menu--;
            if(menu < 0)
                menu = NUM_ITENS - 1;
            desenha_menu(menu);
        }

        if(pressionado(BTN_RB2))
        {
            switch(menu)
            {
                case 0: tela_luz();         break;
                case 1: tela_temperatura(); break;
                case 2: tela_ventilacao();  break;
                case 3: tela_timer();       break;
                case 4: tela_limite();      break;
            }
            desenha_menu(menu);
        }
    }
}
